% !TEX program = xelatex
\documentclass[12pt,a4paper]{ctexart}
\usepackage{amsmath,amssymb,bm}
\usepackage{booktabs}
\usepackage{fancyhdr}
\usepackage{geometry}
\geometry{left=2.5cm,right=2.5cm,top=2.5cm,bottom=2.5cm}

% 页码设置在页脚居中
\pagestyle{fancy}
\fancyhf{}
\fancyfoot[C]{\thepage}
\renewcommand{\headrulewidth}{0pt}
\renewcommand{\footrulewidth}{0pt}

% 公式编号允许在行间使用 \tag
\newcommand{\tagthis}[1]{\tag{#1}}

\begin{document}

\title{\vspace{-2cm}多元线性回归模型的估计——极详细推导与讲解}
\author{}
\date{}
\maketitle

\section{引言}

本文以初学者能跟上的细致程度，讲解多元线性回归模型参数的\textbf{最小二乘估计}、估计量的\textbf{统计性质}（线性性、无偏性、最小方差性）、估计量的\textbf{分布性质}，以及\textbf{随机扰动项方差的估计}。每一步代数变形、每一个矩阵运算、每一个概念转换都给出解释，绝不跳步。

\section{第一步：多元线性回归模型参数的最小二乘估计}

\subsection{1.1 我们要解决什么问题？}

回忆一下简单线性回归：只有一个解释变量 $X$，模型是 $Y_i = \beta_1 + \beta_2 X_i + u_i$，我们用一条直线去拟合数据。

现在情况升级了：我们有\textbf{多个}解释变量 $X_2, X_3, \ldots, X_k$，模型变成了：

\[
Y_i = \beta_1 + \beta_2 X_{2i} + \beta_3 X_{3i} + \cdots + \beta_k X_{ki} + u_i
\]

\begin{itemize}
    \item $Y_i$：被解释变量（因变量）的第 $i$ 个观测值。
    \item $X_{2i}, X_{3i}, \ldots, X_{ki}$：$k-1$ 个解释变量的第 $i$ 个观测值。
    \item $\beta_1$：截距项。
    \item $\beta_2, \beta_3, \ldots, \beta_k$：各个解释变量的偏回归系数。
    \item $u_i$：随机扰动项。
    \item $i = 1, 2, \ldots, n$：共 $n$ 个样本观测值。
\end{itemize}

我们的目标：用样本数据把 $\beta_1, \beta_2, \ldots, \beta_k$ 这些未知参数\textbf{估计出来}。

\textbf{核心思路和简单回归完全一样}：找到一组参数估计值 $\hat{\beta}_1, \hat{\beta}_2, \ldots, \hat{\beta}_k$，使得样本回归函数“尽可能接近”真实数据。而“尽可能接近”的数学标准就是——\textbf{使残差平方和最小}。

\subsection{1.2 残差与残差平方和}

设样本回归函数为：

\[
\hat{Y}_i = \hat{\beta}_1 + \hat{\beta}_2 X_{2i} + \hat{\beta}_3 X_{3i} + \cdots + \hat{\beta}_k X_{ki}
\]

那么第 $i$ 个观测的\textbf{残差}（也叫剩余）就是真实值与估计值之间的差距：

\begin{equation}
e_i = Y_i - \hat{Y}_i = Y_i - (\hat{\beta}_1 + \hat{\beta}_2 X_{2i} + \hat{\beta}_3 X_{3i} + \cdots + \hat{\beta}_k X_{ki}) \tag{3.22}
\end{equation}

\textbf{暂停理解一下}：残差 $e_i$ 就是回归线没“抓住”的那部分。如果回归线完美，$e_i = 0$ 对所有 $i$ 成立。

我们要让\textbf{所有残差的平方加起来}尽可能小：

\[
\sum e_i^2 = \sum [Y_i - (\hat{\beta}_1 + \hat{\beta}_2 X_{2i} + \hat{\beta}_3 X_{3i} + \cdots + \hat{\beta}_k X_{ki})]^2
\]

为什么用平方？和简单回归一样的道理——残差有正有负，直接加会抵消，平方之后全部变正，远的点惩罚更重。

\subsection{1.3 对每个参数求偏导——建立正规方程组}

这是一个关于 $\hat{\beta}_1, \hat{\beta}_2, \ldots, \hat{\beta}_k$ 的多元函数求最小值问题。微积分告诉我们：函数在某点取得极值的\textbf{必要条件}是，对每个变量的偏导数等于零。

所以，对每个 $\hat{\beta}_j$（$j = 1, 2, \ldots, k$）求偏导：

\[
\frac{\partial(\sum e_i^2)}{\partial \hat{\beta}_j} = 0 \quad (j = 1, 2, \ldots, k)
\]

下面我们一个一个来算。

\textbf{对 $\hat{\beta}_1$ 求偏导}：

\[
\frac{\partial}{\partial \hat{\beta}_1} \sum [Y_i - (\hat{\beta}_1 + \hat{\beta}_2 X_{2i} + \cdots + \hat{\beta}_k X_{ki})]^2
\]

方括号内对 $\hat{\beta}_1$ 求导，得到 $-1$（因为 $\hat{\beta}_1$ 的系数是 $-1$），再用链式法则乘以外面的 $2$ 倍方括号本身：

\[
= \sum 2 \cdot [Y_i - (\hat{\beta}_1 + \hat{\beta}_2 X_{2i} + \cdots + \hat{\beta}_k X_{ki})] \cdot (-1) = 0
\]

化简，把 $-2$ 约掉：

\[
-\sum [Y_i - (\hat{\beta}_1 + \hat{\beta}_2 X_{2i} + \cdots + \hat{\beta}_k X_{ki})] = 0
\]

即：

\[
\sum [Y_i - (\hat{\beta}_1 + \hat{\beta}_2 X_{2i} + \cdots + \hat{\beta}_k X_{ki})] = 0
\]

注意方括号里就是 $e_i$，所以这等价于 $\sum e_i = 0$。

\textbf{对 $\hat{\beta}_2$ 求偏导}：

\[
\frac{\partial}{\partial \hat{\beta}_2} \sum [Y_i - (\hat{\beta}_1 + \hat{\beta}_2 X_{2i} + \cdots + \hat{\beta}_k X_{ki})]^2
\]

方括号内对 $\hat{\beta}_2$ 求导，得到 $-X_{2i}$，再用链式法则：

\[
= \sum 2 \cdot [Y_i - (\hat{\beta}_1 + \hat{\beta}_2 X_{2i} + \cdots + \hat{\beta}_k X_{ki})] \cdot (-X_{2i}) = 0
\]

化简：

\[
-\sum X_{2i} [Y_i - (\hat{\beta}_1 + \hat{\beta}_2 X_{2i} + \cdots + \hat{\beta}_k X_{ki})] = 0
\]

同样，方括号里是 $e_i$，所以等价于 $\sum X_{2i} e_i = 0$。

\textbf{对 $\hat{\beta}_k$ 求偏导}（一般情形）：

同理可得：

\[
-\sum X_{ki} [Y_i - (\hat{\beta}_1 + \hat{\beta}_2 X_{2i} + \cdots + \hat{\beta}_k X_{ki})] = 0
\]

等价于 $\sum X_{ki} e_i = 0$。

把所有 $k$ 个方程写在一起：

\[
\begin{cases}
\sum e_i = 0 \\
\sum X_{2i} e_i = 0 \\
\vdots \\
\sum X_{ki} e_i = 0
\end{cases}
\]

\textbf{暂停理解一下}：这 $k$ 个方程叫做\textbf{正规方程组}（normal equations）。它们的含义是——残差与每个解释变量（包括常数项1）的样本内积为零。直观上说，回归线已经把每个解释变量能解释的信息全部“榨取”干净了，剩下的残差跟任何解释变量都不相关。

\subsection{1.4 正规方程组的矩阵表示}

用矩阵来写，可以非常简洁。先把上面的 $k$ 个方程写成向量形式：

\[
\begin{bmatrix}
\sum e_i \\
\sum X_{2i} e_i \\
\vdots \\
\sum X_{ki} e_i
\end{bmatrix}
=
\begin{bmatrix}
0 \\ 0 \\ \vdots \\ 0
\end{bmatrix}
\]

现在看左边这个向量——它的每一行恰好是矩阵 $\mathbf{X}'$ 的一行与残差向量 $\mathbf{e}$ 的乘积。为什么？

回顾数据矩阵 $\mathbf{X}$ 的定义：

\[
\mathbf{X} =
\begin{bmatrix}
1 & X_{21} & \cdots & X_{k1} \\
1 & X_{22} & \cdots & X_{k2} \\
\vdots & \vdots & & \vdots \\
1 & X_{2n} & \cdots & X_{kn}
\end{bmatrix}
\]

$\mathbf{X}$ 是 $n \times k$ 矩阵，第一列全为1（对应截距项），其余列是各解释变量的观测值。

$\mathbf{X}$ 的转置 $\mathbf{X}'$ 是 $k \times n$ 矩阵：

\[
\mathbf{X}' =
\begin{bmatrix}
1 & 1 & \cdots & 1 \\
X_{21} & X_{22} & \cdots & X_{2n} \\
\vdots & \vdots & & \vdots \\
X_{k1} & X_{k2} & \cdots & X_{kn}
\end{bmatrix}
\]

残差向量 $\mathbf{e} = \begin{bmatrix} e_1 \\ e_2 \\ \vdots \\ e_n \end{bmatrix}$ 是 $n \times 1$ 列向量。

那么 $\mathbf{X}'\mathbf{e}$ 就是 $k \times 1$ 向量：

\[
\mathbf{X}'\mathbf{e} =
\begin{bmatrix}
1 \cdot e_1 + 1 \cdot e_2 + \cdots + 1 \cdot e_n \\
X_{21} e_1 + X_{22} e_2 + \cdots + X_{2n} e_n \\
\vdots \\
X_{k1} e_1 + X_{k2} e_2 + \cdots + X_{kn} e_n
\end{bmatrix}
=
\begin{bmatrix}
\sum e_i \\
\sum X_{2i} e_i \\
\vdots \\
\sum X_{ki} e_i
\end{bmatrix}
\]

正好就是左边的向量！所以正规方程组可以极其简洁地写成：

\begin{equation}
\mathbf{X}'\mathbf{e} = \mathbf{0} \tag{3.23}
\end{equation}

\subsection{1.5 从正规方程组解出参数估计量}

现在我们有了 $\mathbf{X}'\mathbf{e} = \mathbf{0}$，但里面含有残差 $\mathbf{e}$，而 $\mathbf{e}$ 又依赖于 $\hat{\boldsymbol{\beta}}$。我们需要把 $\mathbf{e}$ 用 $\hat{\boldsymbol{\beta}}$ 替换掉。

回忆样本回归函数的矩阵形式。总体回归模型的矩阵形式是：

\[
\mathbf{Y} = \mathbf{X}\boldsymbol{\beta} + \mathbf{U}
\]

样本回归函数是：

\[
\hat{\mathbf{Y}} = \mathbf{X}\hat{\boldsymbol{\beta}}
\]

残差向量为：

\[
\mathbf{e} = \mathbf{Y} - \hat{\mathbf{Y}} = \mathbf{Y} - \mathbf{X}\hat{\boldsymbol{\beta}}
\]

对样本回归函数两边同乘 $\mathbf{X}'$：

\begin{equation}
\mathbf{X}'\mathbf{Y} = \mathbf{X}'\mathbf{X}\hat{\boldsymbol{\beta}} + \mathbf{X}'\mathbf{e} \tag{3.24}
\end{equation}

\textbf{暂停理解一下}：这一步就是把 $\mathbf{Y} = \mathbf{X}\hat{\boldsymbol{\beta}} + \mathbf{e}$ 两边左乘 $\mathbf{X}'$，然后用分配律展开。右边出现了 $\mathbf{X}'\mathbf{X}\hat{\boldsymbol{\beta}}$ 和 $\mathbf{X}'\mathbf{e}$ 两项。

由式（3.23），$\mathbf{X}'\mathbf{e} = \mathbf{0}$，所以第二项消失，得到：

\begin{equation}
\mathbf{X}'\mathbf{Y} = \mathbf{X}'\mathbf{X}\hat{\boldsymbol{\beta}} \tag{3.25}
\end{equation}

这就是矩阵形式的\textbf{正规方程组}。

现在要解这个方程，把 $\hat{\boldsymbol{\beta}}$ 解出来。$\mathbf{X}'\mathbf{X}$ 是一个 $k \times k$ 的方阵。由古典假定中的\textbf{无多重共线性假定}，$\mathbf{X}'\mathbf{X}$ 是可逆的，即 $(\mathbf{X}'\mathbf{X})^{-1}$ 存在。

（\textbf{为什么无多重共线性就能保证可逆？} 无多重共线性意味着 $\mathbf{X}$ 的列向量线性无关，而 $\mathbf{X}'\mathbf{X}$ 可逆的充要条件恰好就是 $\mathbf{X}$ 列满秩，即列向量线性无关。）

在式（3.25）两边左乘 $(\mathbf{X}'\mathbf{X})^{-1}$：

\[
(\mathbf{X}'\mathbf{X})^{-1} \mathbf{X}'\mathbf{Y} = (\mathbf{X}'\mathbf{X})^{-1} \mathbf{X}'\mathbf{X}\hat{\boldsymbol{\beta}}
\]

右边 $(\mathbf{X}'\mathbf{X})^{-1} \mathbf{X}'\mathbf{X} = \mathbf{I}_k$（单位矩阵），所以：

\begin{equation}
\hat{\boldsymbol{\beta}} = (\mathbf{X}'\mathbf{X})^{-1}\mathbf{X}'\mathbf{Y} \tag{3.26}
\end{equation}

\textbf{这就是多元线性回归模型参数向量最小二乘估计的矩阵表达式。} 这是一个极其重要、必须牢记的公式。

\textbf{暂停理解一下}：这个公式告诉我们——给定数据 $\mathbf{X}$ 和 $\mathbf{Y}$，只需要做矩阵运算 $(\mathbf{X}'\mathbf{X})^{-1}\mathbf{X}'\mathbf{Y}$，就能一步得到所有参数的OLS估计值。在实际操作中，这就是计算机软件（如EViews、Stata、R）在做的事情。

\subsection{1.6 两个解释变量时的代数表达式}

对于只有两个解释变量的模型 $Y_i = \beta_1 + \beta_2 X_{2i} + \beta_3 X_{3i} + u_i$，$\hat{\beta}_2$ 有一个具体的代数表达式（用离差形式表示）：

\begin{equation}
\hat{\beta}_2 = \frac{\sum y_i x_{2i} \sum x_{3i}^2 - \sum y_i x_{3i} \sum x_{2i} x_{3i}}{\sum x_{2i}^2 \sum x_{3i}^2 - (\sum x_{2i} x_{3i})^2} \tag{3.27}
\end{equation}

其中 $x_{ti} = X_{ti} - \bar{X}_t$，$y_i = Y_i - \bar{Y}$（小写字母表示离差形式）。

\textbf{暂停理解一下}：这个公式看起来复杂，但逻辑是清晰的——当只有一个解释变量时，$\hat{\beta}_2 = \frac{\sum y_i x_i}{\sum x_i^2}$；现在多了一个解释变量 $X_3$，它可能会“混淆” $X_2$ 对 $Y$ 的影响，所以分子和分母都要做“净化”调整——把 $X_2$ 和 $X_3$ 之间的交叉影响考虑进去。分母中的 $\sum x_{2i}^2 \sum x_{3i}^2 - (\sum x_{2i} x_{3i})^2$ 本质上就是在衡量 $X_2$ 和 $X_3$ 之间是否有重叠信息。

\subsection{1.7 例3.1的计算}

书上用2000～2020年中国货币供应量（M2）$Y_t$、GDP $X_{2t}$、CPI $X_{3t}$ 的数据来演示。

把数据排成矩阵：

\[
\mathbf{Y} =
\begin{bmatrix}
134610.3 \\ 158301.9 \\ \vdots \\ 2186795.9
\end{bmatrix},
\quad
\mathbf{X} =
\begin{bmatrix}
1 & 100280.1 & 434.0 \\
1 & 110863.1 & 437.0 \\
\vdots & \vdots & \vdots \\
1 & 1015986.2 & 686.5
\end{bmatrix}
\]

代入公式 $\hat{\boldsymbol{\beta}} = (\mathbf{X}'\mathbf{X})^{-1}\mathbf{X}'\mathbf{Y}$，计算得到：

\[
\hat{\boldsymbol{\beta}} =
\begin{bmatrix}
954083.0 \\ 2.878165 \\ -2601.117
\end{bmatrix}
\]

所以样本回归模型为：

\[
\hat{Y}_i = 954083.0 + 2.878165 X_{2i} - 2601.117 X_{3i}
\]

\textbf{直观理解}：
\begin{itemize}
    \item GDP每增加1亿元，M2平均增加约2.88亿元（CPI不变时）。
    \item CPI每增加1个点，M2平均减少约2601亿元（GDP不变时）——这看起来反直觉，可能是多重共线性等原因导致，后续章节会讨论。
\end{itemize}

\section{第二步：参数最小二乘估计的性质}

从式（3.26）$\hat{\boldsymbol{\beta}} = (\mathbf{X}'\mathbf{X})^{-1}\mathbf{X}'\mathbf{Y}$ 可以看出，参数估计量是样本观测值 $\mathbf{Y}$ 的函数。由于每次抽样的 $\mathbf{Y}$ 不同，$\hat{\boldsymbol{\beta}}$ 也不同——它是随抽样而变化的\textbf{随机变量}。

在古典假定成立的前提下，OLS估计量有三个优良性质：\textbf{线性性、无偏性、最小方差性}。合起来叫做\textbf{最佳线性无偏估计量（BLUE）}。下面逐一证明。

\subsection{2.1 性质一：线性性}

\textbf{目标：证明 $\hat{\boldsymbol{\beta}}$ 是 $\mathbf{Y}$ 的线性函数。}

由式（3.26）：

\[
\hat{\boldsymbol{\beta}} = (\mathbf{X}'\mathbf{X})^{-1}\mathbf{X}'\mathbf{Y}
\]

这里 $(\mathbf{X}'\mathbf{X})^{-1}\mathbf{X}'$ 是一个 $k \times n$ 的矩阵，它\textbf{只依赖于解释变量矩阵 $\mathbf{X}$}。由古典假定，$\mathbf{X}$ 是固定不变的（非随机的）。所以 $(\mathbf{X}'\mathbf{X})^{-1}\mathbf{X}'$ 是一个\textbf{常数矩阵}。

设 $\mathbf{A} = (\mathbf{X}'\mathbf{X})^{-1}\mathbf{X}'$，则 $\hat{\boldsymbol{\beta}} = \mathbf{A}\mathbf{Y}$。

这就意味着 $\hat{\boldsymbol{\beta}}$ 的每一个分量 $\hat{\beta}_j$ 都是 $\mathbf{Y}$ 的各分量 $Y_1, Y_2, \ldots, Y_n$ 的\textbf{线性组合}（加权和），权数仅由 $\mathbf{X}$ 决定。

\textbf{线性性得证。}

\subsection{2.2 性质二：无偏性}

\textbf{目标：证明 $E(\hat{\boldsymbol{\beta}}) = \boldsymbol{\beta}$。}

首先，我们需要一个关键的表达式。把总体模型 $\mathbf{Y} = \mathbf{X}\boldsymbol{\beta} + \mathbf{U}$ 代入 $\hat{\boldsymbol{\beta}}$ 的公式：

\[
\hat{\boldsymbol{\beta}} = (\mathbf{X}'\mathbf{X})^{-1}\mathbf{X}'\mathbf{Y} = (\mathbf{X}'\mathbf{X})^{-1}\mathbf{X}'(\mathbf{X}\boldsymbol{\beta} + \mathbf{U})
\]

展开括号：

\[
= (\mathbf{X}'\mathbf{X})^{-1}\mathbf{X}'\mathbf{X}\boldsymbol{\beta} + (\mathbf{X}'\mathbf{X})^{-1}\mathbf{X}'\mathbf{U}
\]

第一项 $(\mathbf{X}'\mathbf{X})^{-1}\mathbf{X}'\mathbf{X} = \mathbf{I}_k$，所以：

\begin{equation}
\hat{\boldsymbol{\beta}} = \boldsymbol{\beta} + (\mathbf{X}'\mathbf{X})^{-1}\mathbf{X}'\mathbf{U} \tag{3.29}
\end{equation}

\textbf{暂停理解一下}：这个结果非常重要，它是后面推导无偏性和方差的基础。它告诉我们——估计量 $\hat{\boldsymbol{\beta}}$ 等于真实参数 $\boldsymbol{\beta}$ 加上一个由随机扰动 $\mathbf{U}$ 引起的“偏移” $(\mathbf{X}'\mathbf{X})^{-1}\mathbf{X}'\mathbf{U}$。如果 $\mathbf{U} = \mathbf{0}$（没有随机扰动），$\hat{\boldsymbol{\beta}}$ 就精确等于 $\boldsymbol{\beta}$。

现在对两边取期望：

\[
E(\hat{\boldsymbol{\beta}}) = E[\boldsymbol{\beta} + (\mathbf{X}'\mathbf{X})^{-1}\mathbf{X}'\mathbf{U}]
\]

因为 $\boldsymbol{\beta}$ 是常数向量，$(\mathbf{X}'\mathbf{X})^{-1}\mathbf{X}'$ 是常数矩阵（$\mathbf{X}$ 固定），可以把期望算子移进去：

\[
= \boldsymbol{\beta} + (\mathbf{X}'\mathbf{X})^{-1}\mathbf{X}' E(\mathbf{U})
\]

由古典假定中的\textbf{零均值假定}：$E(\mathbf{U}) = \mathbf{0}$，所以：

\begin{equation}
E(\hat{\boldsymbol{\beta}}) = \boldsymbol{\beta} + (\mathbf{X}'\mathbf{X})^{-1}\mathbf{X}' \cdot \mathbf{0} = \boldsymbol{\beta} + \mathbf{0} = \boldsymbol{\beta} \tag{3.30}
\end{equation}

\textbf{无偏性得证。}

\textbf{直观理解}：虽然每一次抽样得到的 $\hat{\boldsymbol{\beta}}$ 不一定等于 $\boldsymbol{\beta}$，但如果我们重复抽样很多很多次，所有这些 $\hat{\boldsymbol{\beta}}$ 的\textbf{平均值}恰好等于真实值 $\boldsymbol{\beta}$。系统上没有偏高也没有偏低的倾向。

\subsection{2.3 性质三：最小方差性（有效性）}

\textbf{目标：证明在所有线性无偏估计量中，OLS估计量 $\hat{\boldsymbol{\beta}}$ 的方差最小。}

这个证明较为繁琐（需要构造任意的线性无偏估计量，然后证明其方差与OLS方差之差是半正定矩阵），书上将其放在附录3.1中。

但结论必须牢记：在古典假定全部满足的条件下，OLS估计量是\textbf{最佳线性无偏估计量（BLUE）}——“最佳”就是指方差最小，“线性无偏”就是前两个性质。

这就是著名的\textbf{高斯-马尔可夫定理}在多元情形下的体现。

\textbf{综合三个性质}：

\begin{center}
\begin{tabular}{lll}
\toprule
\textbf{性质} & \textbf{含义} & \textbf{通俗理解} \\
\midrule
线性性 & $\hat{\boldsymbol{\beta}}$ 是 $\mathbf{Y}$ 的线性组合 & 估计量形式简单 \\
无偏性 & $E(\hat{\boldsymbol{\beta}}) = \boldsymbol{\beta}$ & 平均而言不偏 \\
最小方差性 & 在线性无偏估计中方差最小 & 估计最稳定、最精确 \\
\bottomrule
\end{tabular}
\end{center}

三者合起来 = \textbf{BLUE}（Best Linear Unbiased Estimator）。

\section{第三步：OLS估计的分布性质}

\subsection{3.1 为什么需要知道分布？}

前面我们知道了 $\hat{\boldsymbol{\beta}}$ 的期望和方差，但要进行\textbf{区间估计}（比如算置信区间）和\textbf{假设检验}（比如检验某个系数是否显著），仅有期望和方差是不够的，还需要知道 $\hat{\boldsymbol{\beta}}$ 服从什么\textbf{分布}。

\subsection{3.2 正态性的推导}

由古典假定中的\textbf{正态性假定}，$u_i$ 服从正态分布：$u_i \sim N(0, \sigma^2)$。

因为 $\mathbf{U}$ 的每个分量 $u_i$ 都是正态的，所以：

\[
\mathbf{Y} = \mathbf{X}\boldsymbol{\beta} + \mathbf{U}
\]

$\mathbf{Y}$ 是正态随机向量 $\mathbf{U}$ 的\textbf{线性变换}（$\mathbf{X}\boldsymbol{\beta}$ 是常数，加上 $\mathbf{U}$），所以 $\mathbf{Y}$ 也服从正态分布。

由线性性，$\hat{\boldsymbol{\beta}} = (\mathbf{X}'\mathbf{X})^{-1}\mathbf{X}'\mathbf{Y}$，即 $\hat{\boldsymbol{\beta}}$ 是 $\mathbf{Y}$ 的线性变换。正态分布的线性变换仍然是正态分布。所以 $\hat{\boldsymbol{\beta}}$ 服从正态分布。

\textbf{逻辑链}：$u_i$ 正态 $\Rightarrow$ $\mathbf{U}$ 正态 $\Rightarrow$ $\mathbf{Y}$ 正态 $\Rightarrow$ $\hat{\boldsymbol{\beta}}$ 正态。

\subsection{3.3 方差-协方差矩阵的推导（极其详细，无跳步）}

我们已经知道 $\hat{\boldsymbol{\beta}}$ 服从正态分布，也知道了它的均值 $E(\hat{\boldsymbol{\beta}}) = \boldsymbol{\beta}$。现在要推导它的\textbf{方差-协方差矩阵}。

\textbf{定义回顾}：随机向量 $\mathbf{Z}$ 的方差-协方差矩阵为：

\[
\text{Var-Cov}(\mathbf{Z}) = E\{[\mathbf{Z} - E(\mathbf{Z})][\mathbf{Z} - E(\mathbf{Z})]'\}
\]

对角线元素是各分量的方差，非对角线元素是两两之间的协方差。

\textbf{推导开始}：

由无偏性，$E(\hat{\boldsymbol{\beta}}) = \boldsymbol{\beta}$，所以：

\[
\text{Var-Cov}(\hat{\boldsymbol{\beta}}) = E\{[\hat{\boldsymbol{\beta}} - \boldsymbol{\beta}][\hat{\boldsymbol{\beta}} - \boldsymbol{\beta}]'\}
\]

由式（3.29）：$\hat{\boldsymbol{\beta}} - \boldsymbol{\beta} = (\mathbf{X}'\mathbf{X})^{-1}\mathbf{X}'\mathbf{U}$，代入：

\[
= E\{[(\mathbf{X}'\mathbf{X})^{-1}\mathbf{X}'\mathbf{U}][(\mathbf{X}'\mathbf{X})^{-1}\mathbf{X}'\mathbf{U}]'\}
\]

现在处理转置部分。回忆矩阵转置的性质 $(\mathbf{A}\mathbf{B})' = \mathbf{B}'\mathbf{A}'$：

\[
[(\mathbf{X}'\mathbf{X})^{-1}\mathbf{X}'\mathbf{U}]' = \mathbf{U}'\mathbf{X}[(\mathbf{X}'\mathbf{X})^{-1}]' = \mathbf{U}'\mathbf{X}(\mathbf{X}'\mathbf{X})^{-1}
\]

（这里用到了 $(\mathbf{X}'\mathbf{X})^{-1}$ 是对称矩阵，所以它的转置等于自身。为什么？因为 $(\mathbf{X}'\mathbf{X})' = \mathbf{X}'\mathbf{X}$，即 $\mathbf{X}'\mathbf{X}$ 本身对称，其逆矩阵也对称。）

代入回去：

\[
= E[(\mathbf{X}'\mathbf{X})^{-1}\mathbf{X}'\mathbf{U}\mathbf{U}'\mathbf{X}(\mathbf{X}'\mathbf{X})^{-1}]
\]

现在把常数矩阵移到期望外面。$(\mathbf{X}'\mathbf{X})^{-1}\mathbf{X}'$ 是常数矩阵（在 $\mathbf{X}$ 固定条件下），$(\mathbf{X}'\mathbf{X})^{-1}$ 也是常数矩阵。唯一随机的是 $\mathbf{U}\mathbf{U}'$。所以：

\[
= (\mathbf{X}'\mathbf{X})^{-1}\mathbf{X}' E(\mathbf{U}\mathbf{U}') \mathbf{X}(\mathbf{X}'\mathbf{X})^{-1}
\]

\textbf{关键跳跃点解释}：$E(\mathbf{U}\mathbf{U}')$ 是什么？由古典假定中的\textbf{同方差}和\textbf{无自相关}假定：

\[
E(\mathbf{U}\mathbf{U}') = E
\begin{bmatrix}
u_1^2 & u_1 u_2 & \cdots & u_1 u_n \\
u_2 u_1 & u_2^2 & \cdots & u_2 u_n \\
\vdots & \vdots & \ddots & \vdots \\
u_n u_1 & u_n u_2 & \cdots & u_n^2
\end{bmatrix}
=
\begin{bmatrix}
\sigma^2 & 0 & \cdots & 0 \\
0 & \sigma^2 & \cdots & 0 \\
\vdots & \vdots & \ddots & \vdots \\
0 & 0 & \cdots & \sigma^2
\end{bmatrix}
= \sigma^2 \mathbf{I}_n
\]

\begin{itemize}
    \item 对角线上：$E(u_i^2) = \text{Var}(u_i) = \sigma^2$（同方差假定）
    \item 非对角线上：$E(u_i u_j) = \text{Cov}(u_i, u_j) = 0$（$i \neq j$，无自相关假定）
\end{itemize}

所以 $E(\mathbf{U}\mathbf{U}') = \sigma^2 \mathbf{I}_n$。

代入：

\[
= (\mathbf{X}'\mathbf{X})^{-1}\mathbf{X}' \sigma^2 \mathbf{I}_n \mathbf{X}(\mathbf{X}'\mathbf{X})^{-1}
\]

$\sigma^2 \mathbf{I}_n$ 是标量乘单位矩阵，可以提到最前面；$\mathbf{I}_n \mathbf{X} = \mathbf{X}$（单位矩阵乘任何矩阵等于该矩阵自身）：

\[
= \sigma^2 (\mathbf{X}'\mathbf{X})^{-1}\mathbf{X}'\mathbf{X}(\mathbf{X}'\mathbf{X})^{-1}
\]

$(\mathbf{X}'\mathbf{X})^{-1}\mathbf{X}'\mathbf{X} = \mathbf{I}_k$，所以：

\begin{equation}
\text{Var-Cov}(\hat{\boldsymbol{\beta}}) = \sigma^2 \mathbf{I}_k (\mathbf{X}'\mathbf{X})^{-1} = \sigma^2 (\mathbf{X}'\mathbf{X})^{-1} \tag{3.31}
\end{equation}

\textbf{方差-协方差矩阵推导完毕！}

\[
\text{Var-Cov}(\hat{\boldsymbol{\beta}}) = \sigma^2 (\mathbf{X}'\mathbf{X})^{-1}
\]

\subsection{3.4 单个参数的方差与标准误差}

方差-协方差矩阵的对角线第 $j$ 个元素就是 $\hat{\beta}_j$ 的方差。设 $(\mathbf{X}'\mathbf{X})^{-1}$ 的第 $j$ 行第 $j$ 列元素为 $c_{jj}$，则：

\begin{equation}
\text{Var}(\hat{\beta}_j) = \sigma^2 c_{jj} \tag{3.32}
\end{equation}

\begin{equation}
\text{SE}(\hat{\beta}_j) = \sigma\sqrt{c_{jj}} \tag{3.33}
\end{equation}

其中 $c_{jj}$ 是 $(\mathbf{X}'\mathbf{X})^{-1}$ 中第 $j$ 行第 $j$ 列的元素。

\subsection{3.5 分布结论}

综合以上，在古典假定下，$\hat{\beta}_j$ 服从正态分布：

\begin{equation}
\hat{\beta}_j \sim N[\beta_j, \sigma^2 c_{jj}] \tag{3.34}
\end{equation}

\begin{itemize}
    \item 均值是 $\beta_j$（由无偏性保证）
    \item 方差是 $\sigma^2 c_{jj}$（由方差-协方差矩阵得出）
\end{itemize}

\textbf{暂停理解一下}：知道 $\hat{\beta}_j$ 服从正态分布后，我们就可以对它做标准化变换 $\frac{\hat{\beta}_j - \beta_j}{\sigma\sqrt{c_{jj}}} \sim N(0,1)$，然后用标准正态分布表来进行假设检验和区间估计。但实际上 $\sigma$ 未知，需要用估计值替代，那就涉及 $t$ 分布——这是后续内容。

\section{第四步：随机扰动项方差的估计}

\subsection{4.1 为什么需要估计 $\sigma^2$？}

参数估计量的方差 $\text{Var}(\hat{\beta}_j) = \sigma^2 c_{jj}$ 中，$c_{jj}$ 是可以计算的（给定数据 $\mathbf{X}$ 后，$(\mathbf{X}'\mathbf{X})^{-1}$ 就确定了），但 $\sigma^2$ 是\textbf{总体随机扰动项的方差}，是未知的。

不知道 $\sigma^2$，就没法算出 $\text{Var}(\hat{\beta}_j)$ 的具体数值，也就没法做假设检验和区间估计。所以必须想办法用样本数据来\textbf{估计} $\sigma^2$。

\subsection{4.2 残差向量与残差平方和}

在得到参数估计值 $\hat{\boldsymbol{\beta}}$ 之后，可以计算残差向量：

\[
\mathbf{e} = \mathbf{Y} - \hat{\mathbf{Y}} = \mathbf{Y} - \mathbf{X}\hat{\boldsymbol{\beta}}
\]

残差平方和为：

\[
\sum e_i^2 = \mathbf{e}'\mathbf{e}
\]

\subsection{4.3 $\sigma^2$ 的无偏估计（极其详细推导）}

可以证明（证明见书上附录3.2）：

\[
E\left(\sum e_i^2\right) = E(\mathbf{e}'\mathbf{e}) = (n - k)\sigma^2
\]

\textbf{暂停理解一下}：为什么是 $(n-k)$ 而不是 $n$？因为我们在估计 $k$ 个参数 $\hat{\beta}_1, \hat{\beta}_2, \ldots, \hat{\beta}_k$ 的过程中，消耗了 $k$ 个自由度，所以残差平方和的期望不再是 $n\sigma^2$，而是 $(n-k)\sigma^2$。

由上式可得：

\begin{equation}
E\left(\frac{\sum e_i^2}{n - k}\right) = \sigma^2 \tag{3.35}
\end{equation}

这意味着 $\frac{\sum e_i^2}{n-k}$ 的期望恰好等于 $\sigma^2$——它是一个\textbf{无偏估计量}。

定义：

\begin{equation}
\hat{\sigma}^2 = \frac{\sum e_i^2}{n - k} \tag{3.36}
\end{equation}

则 $\hat{\sigma}^2$ 就是 $\sigma^2$ 的无偏估计。$\hat{\sigma}^2$ 称为\textbf{估计的方差}，$\hat{\sigma}$ 称为\textbf{估计的标准误差}。

\subsection{4.4 参数方差的估计式}

既然 $\sigma^2$ 的无偏估计是 $\hat{\sigma}^2$，我们用 $\hat{\sigma}^2$ 替代 $\sigma^2$，得到参数估计量方差的估计式：

\begin{equation}
\widehat{\text{Var}}(\hat{\beta}_j) = \hat{\sigma}^2 c_{jj} = \left(\frac{\sum e_i^2}{n-k}\right) c_{jj} \tag{3.37}
\end{equation}

\textbf{暂停理解一下}：注意这里多了一个“帽子” $\widehat{\text{Var}}$，表示这是方差的\textbf{估计值}，而不是方差本身。$\text{Var}(\hat{\beta}_j) = \sigma^2 c_{jj}$ 是理论公式（含有未知的 $\sigma^2$），而 $\widehat{\text{Var}}(\hat{\beta}_j) = \hat{\sigma}^2 c_{jj}$ 是可以实际计算的数值。

\subsection{4.5 例3.1的数值计算}

回到货币供应量的例子。$n = 21$（2000～2020年共21个观测值），$k = 3$（截距 + GDP系数 + CPI系数）。

\textbf{步骤1：计算残差平方和}

通过回归计算得到 $\sum e_i^2 = 36{,}623{,}503{,}360$。

\textbf{步骤2：计算 $\hat{\sigma}^2$}

\[
\hat{\sigma}^2 = \frac{\sum e_i^2}{n-k} = \frac{36{,}623{,}503{,}360}{21 - 3} = \frac{36{,}623{,}503{,}360}{18} = 2{,}034{,}639{,}104.00
\]

\textbf{步骤3：计算各参数估计量的方差}

需要知道 $(\mathbf{X}'\mathbf{X})^{-1}$ 对角线上的元素 $c_{11}, c_{22}, c_{33}$。通过计算得到 $c_{22} = 0.00000000002602$（极小，因为GDP数值很大，导致 $\mathbf{X}'\mathbf{X}$ 的元素很大，逆矩阵元素就很小）。

\[
\widehat{\text{Var}}(\hat{\beta}_2) = \hat{\sigma}^2 \cdot c_{22} = 2{,}034{,}639{,}104.00 \times 0.00000000002602 = 0.0529
\]

同理：
\begin{align*}
\widehat{\text{Var}}(\hat{\beta}_1) &= 118{,}064{,}029{,}467.478 \\
\widehat{\text{Var}}(\hat{\beta}_3) &= 682{,}002.861
\end{align*}

\textbf{直观理解}：$\hat{\beta}_2$ 的方差只有0.0529，而 $\hat{\beta}_2 = 2.878165$，标准误约为 $\sqrt{0.0529} \approx 0.23$，相对于系数值2.88来说很小，说明GDP对M2的影响估计得比较精确。而 $\hat{\beta}_1$ 的方差极大（超过千亿元级别），这是因为截距项的估计精度本来就受数据尺度影响很大。

\section{总结}

整篇文章的逻辑链条如下：

\begin{enumerate}
    \item \textbf{估计方法}：用最小二乘法使残差平方和最小，对 $k$ 个参数求偏导得到正规方程组 $\mathbf{X}'\mathbf{e} = \mathbf{0}$，进而解出 $\hat{\boldsymbol{\beta}} = (\mathbf{X}'\mathbf{X})^{-1}\mathbf{X}'\mathbf{Y}$。
    \item \textbf{估计量的性质（BLUE）}：
    \begin{itemize}
        \item \textbf{线性性}：$\hat{\boldsymbol{\beta}}$ 是 $\mathbf{Y}$ 的线性函数，系数只与 $\mathbf{X}$ 有关。
        \item \textbf{无偏性}：$E(\hat{\boldsymbol{\beta}}) = \boldsymbol{\beta}$，利用零均值假定 $E(\mathbf{U}) = \mathbf{0}$ 得到。
        \item \textbf{最小方差性}：在所有线性无偏估计量中方差最小（高斯-马尔可夫定理）。
    \end{itemize}
    \item \textbf{分布性质}：在正态性假定下，$\hat{\beta}_j \sim N[\beta_j, \sigma^2 c_{jj}]$，方差-协方差矩阵为 $\sigma^2(\mathbf{X}'\mathbf{X})^{-1}$。
    \item \textbf{$\sigma^2$ 的估计}：$\hat{\sigma}^2 = \frac{\sum e_i^2}{n-k}$ 是 $\sigma^2$ 的无偏估计，自由度 $n-k$ 反映了估计 $k$ 个参数消耗的自由度。用 $\hat{\sigma}^2$ 替代 $\sigma^2$，就可以算出参数估计量方差的数值。
\end{enumerate}

\end{document}